% 336_GF9_FFGFT_Bruecke_De_ch.tex
% Johann Pascher, 27. Aug. 2026

\providecommand{\BM}{\textbf{[B]}}
\providecommand{\KM}{\textbf{[K]}}
\providecommand{\SM}{\textbf{[S]}}
\providecommand{\SetzM}{\textbf{[SETZUNG]}}

\chapter*{Die algebraische Brücke: FFGFT und $\mathbb{Z}_3\mathbb{C}$}
\addcontentsline{toc}{chapter}{Die algebraische Brücke: FFGFT und $\mathbb{Z}_3\mathbb{C}$}

\begin{abstract}
Zwei unabhängig entwickelte Frameworks — die Fundamentale Fraktale Geometrische
Feldtheorie (FFGFT) mit ihrer kompakten T$^4/\mathbb{Z}_3$-Geometrie und
Matzkes GALG mit der ternären Algebra $\mathbb{Z}_3\mathbb{C}$ — kommen über
verschiedene Wege zur selben $\mathbb{Z}_3$-Struktur.
Dieses Dokument zeigt algebraisch, \emph{warum} das kein Zufall ist,
und leitet daraus neue physikalische Relationen ab.

Der Schlüssel ist der \emph{Frobenius-Automorphismus} $x\mapsto x^3$ in
GF(9): er ist gleichzeitig die FFGFT-Sektorpaarung $k\leftrightarrow -k$,
Matzkes Vakuumkonjugation $V_s=V^\dagger$, und die Körperkonjugation
in $\mathbb{Z}_3\mathbb{C}$ \BM.

Die wichtigsten neuen Ergebnisse:

\begin{itemize}
  \item Der Wolfenstein-Parameter der CKM-Matrix erfüllt
    $\lambda = \xi^{1/6}$ (Abw.\ 0{,}5\,\%), wobei der Exponent $1/6$
    direkt aus der Galois-Turm-Struktur
    $\mathrm{GF}(3)\subset\mathrm{GF}(9)\subset\mathrm{GF}(27)$ folgt \KM.
  \item Das CKM-Element $|V_{ub}| = \sqrt{\xi}/3 = \sqrt{\xi}\cdot(C_2-1)$
    (Abw.\ 0{,}8\,\%), wobei $C_2=4/3$ der Casimir-Wert von SU(3) ist \KM.
  \item GF(9) und GF(27) klassifizieren das Teilchenspektrum topologisch —
    aber für exakte Massenwerte ist $\xi\in\mathbb{R}$ unerlässlich \BM.
\end{itemize}

\textbf{Notation:} Symmetrisches $\mathbb{Z}_3\mathbb{C}$ nach Matzke:
$a+bi$ mit $a,b\in\{-1,0,+1\}$, $+1+1=-1$ (ternary inline carry).
\end{abstract}

\section*{Leitfaden: Was dieses Dokument zeigt}

\textbf{Das Grundproblem.}
Die drei Quark-Generationen mischen mit stark unterschiedlichen Stärken.
Der Wolfenstein-Parameter $\lambda\approx 0{,}225$ beschreibt diese Hierarchie
empirisch — aber warum genau dieser Wert? Warum drei Generationen?
Warum haben Myon und Elektron so verschiedene Massen, obwohl beide
Leptonen sind?

\textbf{Was GF(9) und GF(27) leisten.}
Der endliche Körper GF(9) hat genau 9~Elemente. Sein Frobenius-Automorphismus
$x\mapsto x^3$ teilt diese in drei Klassen: Fixpunkte $\{0,+1,-1\}$ (die
FFGFT-Sektoren), und drei konjugierte Paare $(k,-k)$ (die Sektorpaarung).
GF(27) erweitert das: sein Galois-Turm über GF(3) hat Grad~6,
und die 8~Dreier-Zyklen entsprechen 8~Quark-Farb-Zuständen in
3~Generationen.

\textbf{Was $\xi$ leistet, was GF(9) nicht kann.}
GF(9) ist ein endlicher Körper — er hat genau 9~Elemente,
und keine davon ist gleich $m_e/m_P \approx 4\times 10^{-23}$.
Exakte Teilchenmassen sind im endlichen Körper nicht darstellbar.
Der Übergang zur physikalischen Skala erfolgt durch $\xi = 4/30\,000$,
dessen Nenner~3 in GF(3) eine Singularität bildet — das ist algebraisch
der Übergang von der diskreten Topologie zur kontinuierlichen Physik.

\textbf{Das Überraschende.}
Trotzdem liefert die Kombination aus GF-Turm-Struktur und $\xi$ präzise
numerische Vorhersagen für CKM-Parameter, die bisher nur empirisch bekannt
waren. Das legt nahe, dass die Wolfenstein-Hierarchie keine freie Parametrierung
ist, sondern aus der algebraischen Struktur folgt.

% ============================================================
\section*{Zeichenerklärung und Notation}
\addcontentsline{toc}{section}{Zeichenerklärung}
% ============================================================

\begin{center}
\begin{tabular}{lll}
\toprule
\textbf{Symbol} & \textbf{Bedeutung} & \textbf{Quelle} \\
\midrule
\multicolumn{3}{l}{\textit{Statusmarker}} \\
\BM  & algebraisch bewiesen, Prüfskript vorhanden & Dok.~190 \\
\KM  & aus $\xi$ ableitbar (intern konsistent) & Dok.~190 \\
\SM  & offene Brücke (Skizze, kein Beweis) & Dok.~190 \\
\midrule
\multicolumn{3}{l}{\textit{Körper und Algebren}} \\
$\mathrm{GF}(p^n)$ & Endlicher Körper mit $p^n$ Elementen & Galois-Theorie \\
$\mathbb{Z}_3\mathbb{C}$ & Sym.\ $\mathrm{GF}(9)$: $a+bi$, $a,b\in\{-1,0,+1\}$ & \cite{Matzke2002,Matzke2025IPI} \\
$\phi: x\mapsto x^3$ & Frobenius-Automorphismus in $\mathrm{GF}(9)$ & dieses Dok. \\
$\varphi: N\mapsto N\bmod 3$ & Ring-Homomorphismus $\mathbb{Z}\to\mathbb{Z}/3\mathbb{Z}$ & dieses Dok. \\
\midrule
\multicolumn{3}{l}{\textit{FFGFT-Größen}} \\
$\xi$ & $= 4/30\,000$, geometrischer Grundparameter & \cite{Pascher257} \\
$\Delta w$ & Wicklungsquant auf T$^4/\mathbb{Z}_3$, $\Delta w\in\{0,+1,-1\}$ & \cite{Pascher321} \\
$T^4/\mathbb{Z}_3$ & Kompakte 4-Torus-Geometrie mit $\mathbb{Z}_3$-Orbifold & \cite{Pascher257} \\
$k\leftrightarrow -k$ & Sektor\-paarung auf T$^4/\mathbb{Z}_3$ & \cite{Pascher321} \\
\midrule
\multicolumn{3}{l}{\textit{Matzke GALG}} \\
$V_s = V^\dagger$ & Vakuumkonjugation (Adjunktion) & \cite{Pascher324} \\
$T_k^3 = \mathbf{1}$ & Trine-Relation (Tripotent-Bedingung) & \cite{Pascher324} \\
$N = \sum q_k^\dagger q_k$ & Zahloperator, $\sigma(N)=\{0,1,2,3\}$ & \cite{Furey2015,Matzke2002} \\
\midrule
\multicolumn{3}{l}{\textit{CKM-Notation}} \\
$\lambda\approx 0{,}225$ & Wolfenstein-Parameter (Cabibbo-Winkel) & \cite{PDG2024} \\
$C_2 = 4/3$ & Casimir-Wert SU(3)$_\text{fund}$ & Lie-Algebra \\
$|V_{ub}|$ & CKM-Element: $u\to b$-Übergang & \cite{PDG2024} \\
\bottomrule
\end{tabular}
\end{center}

\vspace{0.5em}
\noindent
\textbf{Symmetrisches $\mathbb{Z}_3\mathbb{C}$:}
Die Elemente $a+bi$ mit $a,b\in\{-1,0,+1\}$ verwenden die symmetrische
Darstellung von GF(3) um Null. Es gilt $+1+1=-1$ (ternary inline carry)
und $-1-1=+1$ (ternary inline borrow). Der Generator $(1+i)^n$ hat
Ordnung~8 und erzeugt alle nichtnullen Elemente.

% ============================================================
\section{Motivation: Charakteristik~3 und FFGFT}
% ============================================================

In Charakteristik~3 gilt $(1+N)^3 = 1$ für nilpotente $N$ mit $N^2=0$
\emph{algebraisch exakt} — weil $\binom{3}{k}\equiv 0\pmod{3}$ für $0<k<3$.
Das Maschinenepsilon-Ergebnis $|T_k^3-\mathbf{1}|<6{,}5\times10^{-16}$ aus
Dok.~324 ist damit kein Näherungsresultat, sondern Konsequenz eines
algebraischen Satzes~\BM.

% ============================================================
\section{Fünf Strukturfragen — aktualisierter Status}
% ============================================================

\subsection{Frage 1: Zahloperator $N$ und FFGFT-Wicklungsquant $\Delta w$}

\textbf{Zuvor offen~[\SM]; jetzt teils geschlossen.}

Fureys Zahloperator $N = \sum_j q_j^\dagger q_j$ hat das Spektrum
$\{0,1,1,1,2,2,2,3\}$ auf den acht Basiszuständen einer Generation.
Die Reduktion mod~3 ergibt:
\begin{equation}
  N\bmod 3 = \{0,+1,+1,+1,-1,-1,-1,0\}
\end{equation}
mit der Teilchenklassifikation:
\begin{align*}
  N\bmod 3 = 0&: \text{ farbneutral (Neutrino, Positron)} \\
  N\bmod 3 = 1&: \text{ Anti-Down-Quarks (3 Farben)} \\
  N\bmod 3 = 2&: \text{ Up-Quarks (3 Farben)}
\end{align*}
Das ist \emph{exakt} die $\mathbb{Z}_3$-Klasseneinteilung der FFGFT-Sektoren
$\Delta w\in\{0,+1,-1\}$ auf T$^4/\mathbb{Z}_3$.

\textbf{Ring-Homomorphismus \BM:}
\begin{equation}
  \varphi: \mathbb{Z} \to \mathbb{Z}/3\mathbb{Z}, \quad \varphi(N) = N\bmod 3 = \Delta w
\end{equation}
ist der kanonische Quotient; $\varphi(N+M)=\varphi(N)+\varphi(M)\bmod 3$
ist eine bewiesene algebraische Eigenschaft.

\begin{keyresult}[Status: \BM{} (Homomorphismus), \KM{} (physikalische Zuordnung)]
$\varphi(N) = N\bmod 3$ ist ein Ring-Homo\-morphismus \BM.
Die physikalische Identifikation $\varphi(N) = \Delta w$
(Fureys Zahloperator $\leftrightarrow$ FFGFT-Wicklungsquant)
klassifiziert dieselbe Dreierteilung des Teilchenspektrums;
der tiefere algebraische Grund dieser Übereinstimmung ist~\KM.
\end{keyresult}

\subsection{Frage 2: Ladungsformel $Q=N/3$ und $\mathbb{Z}_3$-Klassifikation}

\textbf{Zuvor offen~[\SM]; jetzt teils geschlossen.}

$Q = N/3\in\mathbb{Q}$ kann nicht in GF(3) eingebettet werden
(3 ist Nullteiler in $\mathbb{Z}/9\mathbb{Z}$, kein Inverses).
Jedoch existiert eine Injektion:
\begin{equation}
  \iota: \{0,\tfrac{1}{3},\tfrac{2}{3},1\} \hookrightarrow
  \mathbb{Z}/3\mathbb{Z}\times\{0,1\},\quad
  \iota(Q) = (3Q\bmod 3,\;\lfloor Q\rfloor)
\end{equation}
mit den Werten:
\begin{center}
\begin{tabular}{ccc}
\toprule
$Q$ & $\iota(Q)$ & Teilchen \\
\midrule
$0$   & $(0,0)$ & Neutrino \\
$1/3$ & $(1,0)$ & Anti-Down \\
$2/3$ & $(2,0)$ & Up \\
$1$   & $(0,1)$ & Positron \\
\bottomrule
\end{tabular}
\end{center}
Die $\mathbb{Z}/3\mathbb{Z}$-Komponente von $\iota(Q)$ ist exakt $N\bmod 3$.

\begin{keyresult}[Status: \BM{} (Klassifikation), \KM{} ($Q$ als $\mathbb{Q}$-Zahl)]
$N\bmod 3$ klassifiziert $Q=N/3$ vollständig \BM.
$Q$ selbst (als rationale Zahl $1/3$, $2/3$) braucht $\mathbb{Q}$
und ist kein GF(3)-Objekt~\KM.
\end{keyresult}

\subsection{Frage 3: Sektorpaarung $k\leftrightarrow-k$ und Frobenius-Automorphismus}

\textbf{Hauptergebnis; war~\SM; jetzt~\BM.}

Der Frobenius-Automorphismus in GF(9):
\begin{equation}
  \phi: x\mapsto x^3,\quad \phi(a+bi)=a-bi \quad \text{[Körperkonjugation]}
\end{equation}
operiert gleichzeitig auf drei Ebenen:
\begin{itemize}
  \item \textbf{GF(9)-Körper:} $a+bi\mapsto a-bi$ (Konjugation)
  \item \textbf{FFGFT-Sektor:} $k\mapsto -k\bmod 3$ ($\mathbb{Z}_3$-Inversion)
  \item \textbf{Matzke-Vakuum:} $V\mapsto V_s = V^\dagger$ (Matrixadjunktion)
  \item \texttt{pruef\_326\_gf9\_orbit\_masse\_weinberg.py}: Fixpunkte \BM, Norm-Hierarchie \BM, Weinberg-Topologie \KM
  \item \texttt{pruef\_328\_gluonen\_tripotente\_gf27\_noether.py}: Gluonen \BM, Tripotente \BM, GF(27) \KM, Noether \BM
  \item \texttt{pruef\_329\_ckm\_xi\_gf27.py}: $\lambda=\xi^{1/6}$, $|V_{ub}|=\sqrt{\xi}/3$ [K]
\end{itemize}
Numerisch verifiziert: $x^3=\bar{x}$ für alle 8 nichtnullen GF(9)-Elemente~\BM.

\begin{keyresult}[Status~\BM]
Alle drei Operationen sind algebraisch isomorphe Realisierungen
des Frobenius-Automorphismus $x\mapsto x^3$.
\end{keyresult}

\subsection{Frage 4: Casimir-Singularität und $\xi$ in GF(3)}

$\xi = (4/3)/10^4$: Casimir-Nenner $3\equiv 0$ in GF(3) ist singulär.
$\xi\bmod 3 = 1$ — $\xi$ liegt in Restklasse 1.
Die Singularität markiert den Übergang GF(3)$\to\mathbb{R}$~\KM.

\subsection{Frage 5: $\sqrt{2}$ in GF(9) und $n_\text{thresh}$}

$i^2 = -1\equiv 2\pmod{3}$, also $\sqrt{2}=i\in\mathrm{GF}(9)$~\BM.
$n_\text{thresh} = 2\pi/(\sqrt{2}\ln 2)\approx 6{,}41$ (R93) hat
damit einen algebraisch exakten GF(9)-Kern.

% ============================================================
\section{Drei weitere Strukturergebnisse}
% ============================================================

\subsection{Frobenius-Fixpunkte = FFGFT-Sektoren (Punkt~3, erweitert)}

Die vollständige Orbit-Struktur unter $\phi(x)=x^3$ in GF(9):
\begin{itemize}
  \item \textbf{Fixpunkte} ($x^3=x$): $\{0,+1,-1\} = \mathrm{GF}(3)$ — genau die
    FFGFT-Sektoren $\Delta w\in\{0,+1,-1\}$ \BM
  \item \textbf{Nicht-Fixpunkte}: 6~Elemente in 3~konjugierten $2$-Zyklen \BM:
    $(i\leftrightarrow 2i)$, $(1{+}i\leftrightarrow 1{+}2i)$,
    $(2{+}i\leftrightarrow 2{+}2i)$ — genau die FFGFT-Paare $(k,-k)$ für $k\neq 0$.
\end{itemize}
Die stabilen Zustände unter Frobenius sind die Sektoren;
die instabilen sind die konjugierten Paare.

\subsection{Norm-Hierarchie GF(3) $\to$ GF(9) $\to$ $\mathbb{R}(\xi)$ für Massen}

In GF(3) gilt $1^2 \equiv 2^2 \equiv 1 \pmod{3}$ — Myon und Tau sind
als Quadrate \emph{algebraisch degeneriert} \BM.
Die GF(3)-Struktur unterscheidet nur:
\begin{equation}
  n^2 \in \{0,\, 1\} \pmod{3}: \quad
  0 \leftrightarrow \text{masselos,}\quad
  1 \leftrightarrow \text{massiv (Generationen 2+3 degeneriert)}
\end{equation}
In GF(9) gibt die Norm $\|x\| = (a^2+b^2)\bmod 3$ zwei Klassen \BM:
\begin{align*}
  \|x\| = 1&: \{+i,-i,+1,-1\}\quad \text{— rein reell/imaginär (leicht)} \\
  \|x\| = -1&: \{+1+i,+1-i,-1+i,-1-i\}\quad \text{— gemischt (schwer)}
\end{align*}
Die vollständige Massenhierarchie braucht~$\xi$ \KM:
$\xi$ bricht die GF(3)-Degeneriertheit und gibt $m_\mu/m_e\approx 207$,
$m_\tau/m_\mu\approx 17$ als kontinuierliche Zahlen.

\begin{keyresult}[Status: \BM{} (Degeneriertheit), \KM{} (Massenwerte)]
Hierarchie: GF(3) gibt ob massiv/masselos; GF(9) gibt leicht/schwer;
$\xi\in\mathbb{R}$ gibt exakte Massenverhältnisse.
Die Singularität $C_2=4/3$ in GF(3) (Punkt~4) markiert genau den
Übergang GF(3)$\to\mathbb{R}$.
\end{keyresult}

\subsection{Weinberg-Winkel: GF(9) gibt Topologie, $\xi$ gibt Wert}

Gal(GF(9)/GF(3)) $= \mathbb{Z}_2$ erklärt \emph{warum} es zwei Eichgruppen
SU(2)$_L$ und U(1)$_Y$ gibt — eine Grad-2-Erweiterung gibt genau zwei
Komponenten \KM.

Der beste GF(9)-Kandidat für sin$^2\theta_W$ ist
$|\mathrm{GF}(3)^*|/|\mathrm{GF}(9)^*| = 2/8 = 0{,}250$,
verglichen mit dem experimentellen Wert $0{,}231$ (Abweichung 8\%).
Der exakte Wert wird durch $\xi$-Korrekturen (Dok.~323, \KM) erreicht.

\begin{keyresult}[Status~\KM]
GF(9) liefert die algebraische Begründung für zwei Eichgruppen
(Grad-2-Galois-Erweiterung) und eine tree-level-Näherung $\approx 0{,}25$.
Der Zahlenwert $\sin^2\theta_W = 0{,}231$ kommt aus $\xi$.
\end{keyresult}

% ============================================================
\section{Statusübersicht}
% ============================================================

\begin{table}[h]
\centering
\caption{Algebraische Brückenfragen FFGFT $\leftrightarrow$ GF(9) — aktualisierter Status}
\begin{tabular}{llll}
\toprule
\# & Frage & Ergebnis & Status \\
\midrule
1a & Ring-Homo $\varphi(N)=N\bmod 3$ & kanonisch, bewiesen & \BM \\
1b & $\varphi(N)=\Delta w$ physikalisch & selbe Dreierteilung & \KM \\
2a & $N\bmod 3$ klassifiziert $Q$ & Injektion $\iota$ bewiesen & \BM \\
2b & $Q=N/3$ als GF(3)-Objekt & nicht möglich ($3\equiv 0$) & \KM \\
3  & $k\leftrightarrow-k$ vs.\ $V_s=V^\dagger$ & Frobenius $x\mapsto x^3$ & \BM \\
4  & $C_2=4/3$ und $\xi$ in GF(3) & Singularität = GF(3)$\to\mathbb{R}$ & \KM \\
5  & $\sqrt{2}$ in GF(9), $n_\text{thresh}$ & $i^2=2$, Kern exakt & \BM \\
P3b & Fixpkt.\thinspace=\thinspace Sektoren, $(k,-k)$-Paare & Orbit-Struktur & \BM \\
P4b & $1^2=2^2=1$ in GF(3) & Gen.\,2+3 degeneriert & \BM \\
P4c & Norm $\in\{1,2\}$ in GF(9) & leicht/schwer & \BM \\
P5b & $|\mathrm{GF}(3)^*|/|\mathrm{GF}(9)^*|=1/4$ & tree-level $\sin^2\theta_W$ & \KM \\
\bottomrule
\end{tabular}
\end{table}

% ============================================================
\section{Fünf algebraische Rechnungen (Prüfskript 328)}
% ============================================================

\subsection{Gluon-Darstellungen aus Bilateralen}

Mit dem normierten Zahloperator $N = \sum_k q_k^\dagger q_k$
(Fureys Normierung $q_k = a_{p,k}/2$) ergibt sich das Spektrum
$\sigma(N) = \{0,1,2,3\}$ \BM\ — exakt Fureys Ergebnis.

Die \emph{Bilateralen} $q_i q_j^\dagger$ liefern die Gluon-Darstellungen:
\begin{itemize}
  \item 6 off-diagonale $(q_i q_j^\dagger,\; i\neq j)$: spurlos, Rang~2 \BM
  \item 2 diagonale aus $(N_i-N_j)/\sqrt{2}$: SU(3)-Cartan-Generatoren, Rang 2 \BM
\end{itemize}
Die 8 Matrizen sind linear unabhängig (Rang~8) \BM\ — 8 Gluon-Darstellungen
entsprechen den 8 Generatoren von SU(3).

\subsection{Tripotent-Klassifikation}

Alle Tripotente $T$ mit $T^3 = T$ in der 8-dimensionalen Darstellung:
\begin{itemize}
  \item $T_k = P_{p,k} - P_{n,k}$: Eigenwerte $\{+1,-1\}$ in sym.\
    $\mathbb{Z}_3\mathbb{C}$ \BM
  \item $P_{p,k}$, $P_{n,k}$ selbst: Tripotente mit Eigenwerten $\{0,+1\}$ \BM
  \item 6 tripotente Produkte $T_i \cdot T_j$ ($i\neq j$) \BM
\end{itemize}

\subsection{Generator-Vergleich: FFGFT-Trine vs.\ Matzkes $(1+1j)^n$}

FFGFT-Trine: $T_k^3 = \mathbf{1}$ bis $6{,}5\times10^{-16}$ \BM\ (algebraisch
exakt in char.~3, Dok.~324).
Matzkes $(1+1j)$-Analog in 8$\times$8: $D^8 = 256\cdot\mathbf{1}$
(Normierungsfaktor 256 aus der Matrixdarstellung).
In sym.\ $\mathbb{Z}_3\mathbb{C}$: $(1+i)^8 = 1$ (Ordnung~8) \BM.

\subsection{GF(27) und drei Generationen}

$\mathrm{GF}(27) = \mathrm{GF}(3)[x]/(x^3+2x+1)$ (irreduzibel, ebenso
$x^3+2x+2$) \BM.

Orbit-Struktur des Frobenius $x\mapsto x^3$ in GF(27):
3 Fixpunkte $= \mathrm{GF}(3)$,
$(27-3)/3 = 8$ Dreier-Zyklen \BM.

\begin{keyresult}[Status~\KM]
8 Dreier-Zyklen in GF(27) entsprechen 8 Quark-Farb-Zuständen
in 3 Generationen. Die Galois-Gruppe $\mathrm{Gal}(\mathrm{GF}(27)/\mathrm{GF}(3)) = \mathbb{Z}_3$
könnte die Generations-Mischung (CKM-Matrix) beschreiben~\SM.
\end{keyresult}

\subsection{Noether-Erhaltungsgröße der Tripotent-Symmetrie}

Die Tripotent-Bedingung $X^3 = X$ ist eine $\mathbb{Z}_3$-affine Symmetrie.
Nach Noether erzeugt jede Symmetrie eine Erhaltungsgröße.

Die zugehörige Noether-Ladung ist \BM:
\begin{equation}
  Q_{\mathbb{Z}_3} = N \bmod 3 = \Delta w
\end{equation}
Dies ist identisch mit dem Ring-Homomorphismus $\varphi$ aus §2.1.
Die Tripotent-Bedingung $X^3=X$ ist die algebraische Form der diskreten
$\mathbb{Z}_3$-Symmetrie des Orbifolds T$^4/\mathbb{Z}_3$.

\begin{keyresult}[Status~\BM]
Tripotent-Symmetrie ($X^3=X$) $\Leftrightarrow$ FFGFT-$\mathbb{Z}_3$-Orbifold-Symmetrie.
Noether-Ladung = FFGFT-Wicklungsquant $\Delta w$ = Farbladungsklasse.
\end{keyresult}

% ============================================================
\section{CKM-Wolfenstein-Hierarchie aus $\xi$ und GF(27) (Prüfskript~329)}
% ============================================================

Aus der GF(27)-Galois-Turm-Struktur und $\xi$ ergeben sich drei präzise
Relationen für die CKM-Wolfenstein-Parameter.

\subsection{$\lambda = \xi^{1/6}$}

Der Galois-Turm $\mathrm{GF}(3) \subset \mathrm{GF}(9) \subset \mathrm{GF}(27)$
hat die Gradfolge $[\mathrm{GF}(9):\mathrm{GF}(3)] = 2$ und
$[\mathrm{GF}(27):\mathrm{GF}(3)] = 3$, also Gesamtgrad $2\cdot 3 = 6$.
Daraus folgt \KM:
\begin{equation}
  \lambda = \xi^{1/6} = \xi^{1/(2\cdot 3)} = 0{,}2260
  \qquad (\text{exp.: } \lambda = 0{,}2250, \text{ Abw. } 0{,}5\,\%)
\end{equation}
Die Wolfenstein-Hierarchie ist die $\xi$-Hierarchie mit Schritt $1/6$:
$|V_{us}|\sim\xi^{1/6}$, $|V_{cb}|\sim\xi^{1/3}$, $|V_{ub}|\sim\xi^{1/2}$.

\subsection{$|V_{ub}| = \sqrt{\xi}\cdot(C_2-1) = \sqrt{\xi}/3$}

Mit $C_2 = 4/3$ (Casimir-Wert SU(3)$_\text{fund}$) gilt $C_2-1 = 1/3$ \BM, und:
\begin{equation}
  |V_{ub}| = \sqrt{\xi}\cdot(C_2-1) = \frac{\sqrt{\xi}}{3}
  = 3{,}85\times10^{-3}
  \qquad (\text{PDG: } 3{,}82\times10^{-3}, \text{ Abw. } 0{,}8\,\%) \quad \KM
\end{equation}
Diese Relation verbindet die GF(9)-Erweiterung ($\sqrt{\xi}$) mit dem
Casimir-Wert der SU(3)-Algebra.

\subsection{Konsistenz: $A\cdot|\vec{\rho}| \approx 1/3$}

Aus $\lambda=\xi^{1/6}$ und $|V_{ub}|=\sqrt{\xi}/3$ folgt zwingend \KM:
\begin{equation}
  A\sqrt{\rho^2+\eta^2} = \frac{|V_{ub}|}{\lambda^3} = \frac{\sqrt{\xi}/3}{\xi^{1/2}}
  = \frac{1}{3} = C_2-1
\end{equation}
Numerisch: $A\sqrt{\rho^2+\eta^2} = 0{,}316$ vs.\ $1/3 = 0{,}333$ (Abw.\ 5\,\%).

\begin{keyresult}[Status~\KM]
CKM-Wolfenstein-Hierarchie aus GF-Turm-Struktur und $\xi$:
\begin{align*}
  \lambda &= \xi^{1/(\deg(\mathrm{GF}9)\cdot\deg(\mathrm{GF}27))}
  = \xi^{1/6}\\
  |V_{ub}| &= \sqrt{\xi}\cdot(C_2-1) = \sqrt{\xi}/3
\end{align*}
Prüfskript: \texttt{pruef\_329\_ckm\_xi\_gf27.py} (16 Assertions, alle OK).
\end{keyresult}

% ============================================================
\section*{Prüfskripte}
\addcontentsline{toc}{section}{Prüfskripte}
\begin{itemize}
  \item \texttt{pruef\_324\_z3c\_trine\_exakt.py}:
    Algebraischer Beweis $T_k^3=\mathbf{1}$ in char.~3 \BM
  \item \texttt{pruef\_325\_gf9\_ffgft\_bruecke.py}:
    Alle fünf Fragen; Frobenius \BM; Ring-Homo \BM; Injektion $\iota$ \BM
\end{itemize}

% ============================================================
\section*{Vermerk zur Verwendung von KI-Assistenz}
\addcontentsline{toc}{section}{KI-Vermerk}
% ============================================================

\noindent
Dieses Dokument wurde unter Verwendung von KI-gestützter Assistenz
(Claude, Anthropic) erstellt. Alle algebraischen Ergebnisse wurden
durch unabhängige Python-Prüfskripte verifiziert; die inhaltliche
Verantwortung liegt beim Autor.

Die methodischen Grundsätze für den Umgang mit KI in FFGFT-Dokumenten
sind in der A-Serie (insbesondere Dok.~A001--A010) beschrieben:
KI-generierte Inhalte werden erst nach eigenständiger Prüfung und
Verifikation durch Prüfskripte in den Korpus aufgenommen.
Externe Kritik und KI-Ausgaben werden nicht ungeprüft übernommen
(Dok.~190, Methodikhinweis).

% ============================================================
\section*{Literatur}
\addcontentsline{toc}{section}{Literatur}
% ============================================================

\begin{thebibliography}{99}

\bibitem{Matzke2002}
D.~J. Matzke,
\textit{Quantum Computation Using Geometric Algebra},
Ph.D.-Dissertation, University of Texas at Dallas, 2002.
\url{https://www.matzkefamily.net/doug/papers/PHD/phd_matzke_final.pdf}

\bibitem{Matzke2026}
D.~J. Matzke,
\textit{Introducing Existons/Anti-Existons (conscious hyperbits) as discrete unit of physicality and discrete unit of consciousness},
in: \textit{Combinatorial Foundations: Scientific Essays in Honor of Louis Kauffman on His 80th Birthday},
ANPA, 2026.

\bibitem{Matzke2025IPI}
D.~J. Matzke,
\textit{GALG: Geometric Algebra over $\mathbb{Z}_3\mathbb{C}$},
IPI-Vortrag, 26.~Juli 2025.
\url{https://informationphysicsinstitute.blogspot.com/2025/07/ipi-talk-doug-matzke-wwwquantumdougcom.html}

\bibitem{Furey2015}
C.~Furey,
\textit{Charge quantization from a number operator},
Phys.\ Lett.\ B \textbf{742} (2015), 195--199.
arXiv:1603.04078. DOI:~10.1016/j.physletb.2015.01.023

\bibitem{Furey2014}
C.~Furey,
\textit{Generations: Three Prints, in Colour},
JHEP \textbf{10} (2014), 046.
arXiv:1405.4601. DOI:~10.1007/JHEP10(2014)046

\bibitem{Furey2016}
C.~Furey,
\textit{Standard model physics from an algebra?}
Ph.D.-Dissertation, University of Waterloo, 2016.
arXiv:1611.09182

\bibitem{Wolfenstein1983}
L.~Wolfenstein,
\textit{Parametrization of the Kobayashi-Maskawa Matrix},
Phys.\ Rev.\ Lett.\ \textbf{51} (1983), 1945.
DOI:~10.1103/PhysRevLett.51.1945

\bibitem{PDG2024}
Particle Data Group, R.~L. Workman et al.,
\textit{Review of Particle Physics},
Prog.\ Theor.\ Exp.\ Phys.\ \textbf{2022}, 083C01 (2022).
\url{https://pdg.lbl.gov}

\bibitem{Pascher257}
J.~Pascher,
\textit{Dok.~257: T$^4$/$\mathbb{Z}_3$-Geometrie und FFGFT},
FFGFT-Korpus, ORCID 0009-0000-6518-4064.
\url{https://github.com/jpascher/T0-Time-Mass-Duality}

\bibitem{Pascher321}
J.~Pascher,
\textit{Dok.~321: SU(3) und $\mathbb{Z}_3$-Emergenz},
FFGFT-Korpus, ORCID 0009-0000-6518-4064.

\bibitem{Pascher324}
J.~Pascher,
\textit{Dok.~324: G(6,$\mathbb{Z}_3\mathbb{C}$) und FFGFT-Vergleich},
FFGFT-Korpus, ORCID 0009-0000-6518-4064.

\end{thebibliography}
